\setcounter{section}{7}

\Needspace{5\baselineskip}
\hypertarget{sequence-to-sequence-seq2seq-learning}{%
\section{Sequence-to-Sequence (Seq2Seq) Learning}\label{sequence-to-sequence-seq2seq-learning}}

Technically, an encoder converts a variable-length input sequence into a fixed-shape context variable \(\mathbf{c}\), encoding the input's information in that variable. An RNN can serve as the encoder.

Consider a sample consisting of one sequence (batch size 1). Let the input sequence be \(x_1,\ldots,x_T\), where \(x_t\) is its \(t\)th token. At time \(t\), the RNN transforms token \(x_t\)'s input feature vector \(\mathbf{x}_t\) and \(\mathbf{h}_{t-1}\) (the preceding hidden state) into \(\mathbf{h}_t\) (the current hidden state). Let function \(f\) describe this recurrent-layer transformation:

\[
\mathbf{h}_t = f(\mathbf{x}_t,\mathbf{h}_{t-1})
\]

\Needspace{5\baselineskip}
In summary, the encoder uses a chosen function \(q\) to transform hidden states from all time steps into the context variable:

\[
\mathbf{c} = q(\mathbf{h}_1,\ldots,\mathbf{h}_T)
\]

As noted above, encoder output \(\mathbf{c}\) encodes the entire input sequence \(x_1,\ldots,x_T\). For a training output sequence \(y_1,y_2,\ldots,y_{T'}\), at each time \(t'\) (distinct from input/encoder time \(t\)), the probability of decoder output \(y_{t'}\) depends on preceding outputs \(y_1,\ldots,y_{t'-1}\) and context \(\mathbf{c}\), namely \(P(y_{t'} \mid y_1,\ldots,y_{t'-1},\mathbf{c})\).

To model this conditional probability over a sequence, use another RNN as the decoder. At any output time \(t'\), it receives the preceding output \(y_{t'-1}\) and context \(\mathbf{c}\), transforming these with preceding hidden state \(\mathbf{s}_{t'-1}\) into current state \(\mathbf{s}_{t'}\). Function \(g\) can therefore represent the decoder hidden-layer transformation:

\[
\mathbf{s}_{t'} = g(y_{t'-1},\mathbf{c},\mathbf{s}_{t'-1})
\]

After obtaining the decoder hidden state, an output layer and softmax compute the conditional distribution \(P(y_{t'} \mid y_1,\ldots,y_{t'-1},\mathbf{c})\) of output \(y_{t'}\) at time \(t'\).

At every step, the decoder predicts an output-token distribution. As in language modeling, softmax provides the distribution and cross-entropy loss is used for optimization. To load different-length sequences into identically shaped mini-batches, special padding tokens are appended to their ends. Their predictions should be excluded from the loss calculation. The sequence\_mask function masks irrelevant entries by setting them to zero, so subsequent computations for irrelevant predictions multiply by zero and yield zero (cross-entropy loss -sigma(pi*logqi), pi=0).

\FloatBarrier
\Needspace{5\baselineskip}
\hypertarget{references-24}{%
\subsection{References}\label{references-24}}

\vspace{0.25em}
\begin{itemize}
\tightlist
\item
  Zhang, A., Lipton, Z. C., Li, M., \& Smola, A. J. (2023). \href{https://D2L.ai}{Dive into Deep Learning}. Cambridge University Press.
\item
  Sutskever, I., Vinyals, O., \& Le, Q. V. (2014). \href{https://proceedings.neurips.cc/paper_files/paper/2014/hash/5a18e133cbf9f257297f410bb7eca942-Abstract.html}{Sequence to Sequence Learning with Neural Networks}. NeurIPS 2014.
\end{itemize}

\clearpage
